\section{The curve}\label{sec:curve}

\begin{proposition}[\cid{A.1}; exact computation]\label{prop:curve}
The following data hold for $E\colon y^2=x^3-22735x^2+106272484x$.
\begin{enumerate}[label=(\roman*)]
\item \emph{Minimal model.} The substitution $x=X+7578$ (with $u=1$, so the invariant differential is unchanged) gives the global minimal model
\[
y^2=X^3-X^2-66020924X-65074271436 .
\]
Since $X=x-7578$, the basis points $P_1=(4547,327384)$ and $P_2=(5184,281520)$ of Theorem~\ref{thm:mw} have coordinates $(-3031,327384)$ and $(-2394,281520)$ on this model. Its invariants are $c_4=3169004368$, $c_6=56243184546880$ and
\[
\Delta=2^8\cdot3^2\cdot2243\cdot4547^3\cdot5843^2=16586638656786071143743744>0 .
\]
\item \emph{Bad reduction.} The bad primes behave as follows.
\begin{center}
\begin{tabular}{@{}lllll@{}}
\toprule
$\ell$ & Kodaira type & $c_\ell$ & $a_\ell$ & remark\\ \midrule
$2$ & $\mathrm{I}_1^*$ & $2$ & $0$ & additive\\
$3$ & $\mathrm{I}_2$ & $2$ & $-1$ & nonsplit multiplicative\\
$2243$ & $\mathrm{I}_1$ & $1$ & $+1$ & split multiplicative\\
$4547$ & $\mathrm{III}$ & $2$ & $0$ & additive, potentially good supersingular\\
$5843$ & $\mathrm{I}_2$ & $2$ & $-1$ & nonsplit multiplicative\\ \bottomrule
\end{tabular}
\end{center}
The Tamagawa product is $16$. At $4547$ one has $j\equiv1728$, $4547\equiv3\pmod 4$, and inertia acts through a cyclic group of order $4$.
\item \emph{Global invariants.} $E(\Q)_{\mathrm{tors}}=\Z/2$, generated by $T=(0,0)$. The root number is $+1$.
\item \emph{Good primes.} At every good prime $p$ the trace $a_p$ is even, and $p\nmid\#E(\F_p)$. The good supersingular primes below $2000$ are
\[
23,\,29,\,131,\,167,\,281,\,563,\,661,\,823,\,839,\,863,\,887,\,1013,\,1471,\,1511,\,1823 .
\]
\end{enumerate}
\end{proposition}

\begin{proof}[Justification]
Items (i) and (iv) (for $p<2000$) are recomputed in exact integer arithmetic by \path{tools/replay/curve_invariants.py}, which uses no computer algebra system. The Kodaira types, Tamagawa numbers and root number are the output of Tate's algorithm and the local root-number formulas, recorded in \path{certificate/hypotheses_output.txt}.

For (iv) in general: the rational $2$-torsion point makes $\#E(\F_p)$ even, so $a_p=p+1-\#E(\F_p)$ is even. If $p\mid\#E(\F_p)$, then $\#E(\F_p)$ is a positive even multiple of $p$ in $[p+1-2\sqrt p,\,p+1+2\sqrt p]$. That forces $2p\le p+1+2\sqrt p$, hence $p\le5$; and $\#E(\F_5)=8$.
\end{proof}

\begin{proposition}[\cid{A.2}; exact computation]\label{prop:noncm}
$j(E)=1322372123614064/689197281363\notin\Z$, so $E$ has no complex multiplication.
\end{proposition}

\begin{theorem}[\cid{A.3}; written proof applying published theorems]\label{thm:image}
For every odd prime $\ell$, the mod-$\ell$ representation $\bar\rho_{E,\ell}\colon\Gal(\bar\Q/\Q)\to\GL_2(\F_\ell)$ is surjective. Hence the $\ell$-adic representation is surjective onto $\GL_2(\Z_\ell)$ for $\ell\ge5$.
\end{theorem}
\begin{proof}
Reducibility would give a rational cyclic isogeny of odd prime degree $\ell$. The prime-isogeny theorem restricts $\ell$ to $3,5,7,11,13,17,19,37,43,67,163$. Direct point counts give the following witnesses:
\begin{center}\begin{tabular}{@{}llrr@{}}\toprule
$\ell$ & $q$ & $a_q$ & $a_q^2-4q$\\\midrule
$3,7,11,19,43,67,163$ & $5$ & $-2$ & $-16$\\
$13,17,37$ & $7$ & $-2$ & $-24$\\
$5$ & $17$ & $-6$ & $-32$\\\bottomrule
\end{tabular}\end{center}
Every displayed discriminant is a nonsquare modulo each listed $\ell$, and $q\ne\ell$. Thus the indicated Frobenius has no invariant line, contradicting reducibility. These finite checks are reproduced in \path{review/pro_report/independent_checks.py}.

For $\ell\ne2243$, multiplicative reduction with $v_{2243}(\Delta)=1$ supplies a nontrivial Tate transvection. For $\ell=2243$, use instead $5843$, with $v_{5843}(\Delta)=2$ prime to $\ell$; nonsplit multiplicative reduction is split over an unramified quadratic extension, so its inertia gives the same conclusion. An irreducible subgroup containing a transvection also contains a conjugate fixing a different line. In a basis of those two lines their powers give all upper and lower elementary unipotents over the prime field $\F_\ell$. These generate $\mathrm{SL}_2(\F_\ell)$. The determinant is the surjective cyclotomic character, proving the mod-$\ell$ assertion. The standard integral lifting theorem for closed subgroups at $\ell\ge5$, with the cyclotomic determinant, gives the $\ell$-adic assertion. These are the retained rational-isogeny and lifting imports of \cid{A.3}; no lift is asserted at $\ell=3$.
\end{proof}

\begin{theorem}[\cid{A.4}; written proof applying published theorems]\label{thm:manin}
The optimal curve in the isogeny class of $E$ has Manin constant $c_0=2^a$ with $a\le 8$; in particular $4547\nmid c_0$. The chosen parametrization $\phi_E$ specified in Appendix~\ref{app:latticeproof} has Manin constant $c_E$ dividing $2^9$. Consequently, with
\[
\lambda(r)=2\pi i\int_{i\infty}^{r}f(z)\,dz,\qquad [r]^{+}=\frac{\mathrm{Re}\,\lambda(r)}{\Omega^{+}},\qquad [r]^{-}=\frac{\mathrm{Im}\,\lambda(r)}{\Omega^{-}}\qquad(r=a/d\in\Q,\ (a,d)=1),
\]
where $\Omega^+=\omega_1$ and $\Omega^-$ are the least positive real and imaginary periods, $[r]^\pm\in\frac1{1024}\Z$ whenever $\gcd(d,N)$ is an exact divisor of $N$. These are the rational unitary cusps needed by the retained certificates, including denominators $p^n$ for $p\nmid N$ or $p\,\|\,N$.
\end{theorem}
The complete derived argument, with its geometric imports and exact chosen maps, is in Appendix~\ref{app:latticeproof}. The bound $a\le8$ is a derived lemma; no single source states it. It combines the method of Pasten \cite[\S10]{pasten} with \cite[6.12(c)]{cns} and \cite[3.6(a)]{ars}. The exclusion $4547\nmid c_0$ also follows from Edixhoven's thesis \cite[Thm.~4.6.3]{edixhoventhesis}, whose statement was checked directly in the follow-up review; the published version \cite[Thm.~3]{edixhoven} was read through restatements \cite{cns}. The stronger bound $a\le4$ would need an unverified contact-order input and is not used.

For these cusps, Manin--Drinfeld and rationality put their images in $E(\Q)_{\rm tors}$, of exponent $2$. Therefore $2c_E\lambda(r)$ is a period. The rectangular period lattice ($\Delta>0$) gives the stated bound separately for the real and imaginary parts. This argument makes no assertion about all nonrational cusps.

\begin{remark}[Why the lattice matters]\label{rem:lattice}
The identification step needs both a proved lattice and an interval containing exactly one point of that lattice. Radius $<1/2048$ supplies uniqueness, not membership. An incorrect denominator, whether a power of $2$ or an odd factor, can silently produce a wrong rounded value. For example an interval centered at $0$ of radius $2^{-19}$ contains the true value $2^{-20}$ but has $0$ as its only point of $\frac1{1024}\Z$. Thus the Manin-constant and cusp arguments remain load-bearing even when all rounding checks pass.
\end{remark}
